% =====================================================================
%  Is the Riemann Hypothesis Solved?
%  Front matter: half-title, title page, copyright, dedication, epigraph,
%  contents, figures, tables, preface, how to use.
% =====================================================================

% ---- Half-title ---------------------------------------------------------
\pagestyle{plain}
\clearpage
\thispagestyle{empty}
\vspace*{1.15in}
{\sffamily\fontsize{22}{26}\selectfont\color{txt}\bookTitle\par}
\vspace{10pt}
{\sffamily\fontsize{11}{14}\selectfont\color{txtdim}\bookSubtitle\par}
\vfill
{\sffamily\small\color{txtdim}\bookAuthor\par}
\vspace{6pt}
{\sffamily\small\color{txtdim}\bookImprint\par}
\clearpage

% ---- Title page ---------------------------------------------------------
\clearpage
\thispagestyle{empty}
\vspace*{0.75in}
{\sffamily\fontsize{26}{30}\selectfont\bfseries\color{txt}\bookTitle\par}
\vspace{12pt}
{\sffamily\fontsize{12}{16}\selectfont\color{accent}\bookSubtitle\par}
\vspace{22pt}
\noindent{\color{accent}\rule{\textwidth}{1.8pt}}\par
\vspace{18pt}
{\sffamily\large\color{txt}\bookAuthor\par}
\vspace{8pt}
{\sffamily\color{txtdim}\bookImprint, publisher\par}
\vspace{4pt}
{\sffamily\color{txtdim}\bookEdition\par}
\vspace{18pt}
{\sffamily\small\color{txtdim}
Facts current as of \bookFactsDate.\par}
\vfill
\clearpage

% ---- Copyright page -----------------------------------------------------
\clearpage
\thispagestyle{plain}
\noindent{\sffamily\bfseries\color{txt}\bookTitle}\par
\noindent{\sffamily\small\color{txtdim}\bookSubtitle}\par
\vspace{14pt}
\noindent\rule{\textwidth}{0.4pt}\par
\vspace{10pt}
\noindent Copyright \textcopyright\ \bookCopyright\ \bookAuthor. All rights reserved.\par
\vspace{6pt}
\noindent \bookImprint, publisher.\par
\vspace{10pt}
\noindent\textbf{Edition:} \bookEdition.\par
\vspace{6pt}
\noindent\textbf{Subject headings:} Riemann hypothesis --- Mathematical
verification; Number theory --- Computational methods; Formalization
(Computer science); Machine learning --- Mathematical applications.\par
\vspace{10pt}
\noindent\textbf{Licence.} This work is published under the Apache License,
Version 2.0. The complete licence text is available at
\href{https://www.apache.org/licenses/LICENSE-2.0}{\texttt{apache.org/licenses/LICENSE-2.0}}.
The book may be copied, modified and redistributed, including for commercial
purposes, provided that this notice and the attribution required by the licence
are retained.\par
\vspace{10pt}
\noindent\textbf{DOI:} \href{https://doi.org/\bookDOI}{\texttt{\bookDOI}}\par
\vspace{10pt}
\noindent\textbf{ISBN:} 978-1-0000-0000-0 (placeholder; to be assigned on
publication).\par
\vspace{10pt}
\noindent\textbf{Source code.} No source repository is published for this
edition. The Lean development examined in this book is the work of its author and
is not redistributed here; it is cited by the author's own distribution point and
by the third-party mirror recorded in the Notes and References.\par
\vspace{10pt}
\noindent\textbf{What this book is.} This book reports and analyses a claim. It
does not assert that the Riemann Hypothesis has been proved, and it does not assert
that the claim is false. Chapter~13 states a finding; the finding is that the
matter is unresolved on the available evidence, and the book says precisely what
would settle it. Quotations from the document under examination are reproduced for
the purposes of criticism and review. Every fact concerning the mathematics is
taken from a public or peer-reviewed source, identified in the Notes and
References.\par
\vspace{10pt}
\noindent\textbf{Disclaimers.} The names of mathematicians, institutions and
software projects mentioned in this book are the names of the people and things
whose work is being discussed. Their inclusion is not a statement that they
endorse, support, or have verified anything described here. In particular, the
author of the claim under examination is thanked in that claim for a list of
advisers; this book reports that list as part of the document, and draws no
inference from it.\par
\vspace{10pt}
\noindent\textbf{Facts current as of} \bookFactsDate. Statements that are likely to
change --- the status of independent compilation, the status of independent
review, the state of the underlying literature --- are date-stamped where they
appear. Nothing in this book should be read as a statement about events after
that date.\par
\vspace{8pt}
\noindent{\color{accent}\rule{\textwidth}{0.4pt}}\par
\vspace{8pt}
{\small\color{txtdim}\sffamily
\noindent A note on the subject's provenance. The document examined in this book
is the work of Levent Alpöge, who published it under his own name and made it
available for public download. This book discusses it in order to evaluate it.
It reproduces short quotations under the fair-use provisions for criticism and
review, attributes every quotation, and links to the original in every case
where the source is a document.\par}

% ---- Dedication ---------------------------------------------------------
\clearpage
\thispagestyle{empty}
\vspace*{2.5in}
  \begin{center}
    {\large\itshape\color{txt}
      To the people who read the code instead of the headline,\\[4pt]
      and to the ones who do both.\par}
  \end{center}
  \vfill
\clearpage

% ---- Epigraph ------------------------------------------------------------
\clearpage
\thispagestyle{empty}
\vspace*{2.0in}
  \begin{center}
    \begin{minipage}{0.88\textwidth}
      \raggedright
      {\sffamily\color{accent}\rule{\textwidth}{2pt}}\par
      \vspace{14pt}
      {\large\itshape\color{txt}
      ``I am presenting the world with a complete argument before having mastered
      a massive neighbourhood around its mathematics.''\par}
      \vspace{12pt}
      {\raggedleft\sffamily\small\color{txtdim}
      Levent Alpöge, \emph{Remark}, 2026\par}
    \end{minipage}
  \end{center}
  \vfill
\clearpage

% ---- Contents -----------------------------------------------------------
\clearpage
\pdfbookmark[0]{Contents}{contents}
\tableofcontents

% ---- List of figures ----------------------------------------------------
\clearpage
\pdfbookmark[0]{List of Figures}{list-of-figures}
\listoffigures

% ---- List of tables -----------------------------------------------------
\clearpage
\pdfbookmark[0]{List of Tables}{list-of-tables}
\listoftables

% ---- Preface ------------------------------------------------------------
\chapter*{Preface}
\addcontentsline{toc}{chapter}{Preface}
\markboth{Preface}{Preface}

\thispagestyle{plain}

On 20 November 2026 a mathematician working at Anthropic published a four-page
note claiming to have proved the Riemann Hypothesis, together with a 239-megabyte
Lean development containing 175{,}187 source files. Within hours, two
university-based mathematicians had compiled it. Within days, the world's
newspapers had called it the greatest mathematical result of the century, and
the remark's own author was describing it as ``one dawn across science of
something far greater''.

That combination --- an astonishing claim, a real machine-checkable artifact,
instant global attention, and no peer review --- is rare, and it is the reason
this book exists. It is not a book about the Riemann Hypothesis. There are
several of those, and they are better at being about the Riemann Hypothesis.
It is a book about \emph{what it takes to know whether a machine-checked proof
of the Riemann Hypothesis is a proof}, and about a specific case in which the
question can be asked concretely.

Three commitments govern how the book is written.

The first is \textbf{that the artifact is taken seriously}. Most commentary on
claims like this one consists of either applause or dismissal. Both are cheap.
The Lean development in question is unusually substantial: it contains no axioms
and no admitted gaps, it is organised around genuine mathematics, and its
top-level statement is Mathlib's own definition of the Riemann Hypothesis
rather than a convenient restatement of something weaker. That is worth
something. A book that treated it as obviously fraudulent would be no more
useful than one that treated it as obviously right.

The second is \textbf{that the code is read, not summarised}. Every Lean
statement quoted in this book was read in the submitted source tree, at the
pinned revision, and the dependency chain from the headline theorem down to the
classical fixed-point formula was traced by hand and reported in full, in
Chapter~7. This is the part of the work that is reproducible by any reader with
the same tools.

The third is \textbf{that the question is answered without flinch}. The answer,
reached in Chapter~13, is that the claim is \emph{unresolved}. The artifact
compiles; nobody has audited the mathematics; and the author's own account of
the method records an a-priori ceiling at two thirds of the zeros, and an
obstruction whose stated remedy requires a conjecture that is still open. That is
not a verdict of ``false''. It is a verdict of ``not established'', accompanied
by a protocol --- Chapter~14 --- that would settle it.

Readers who want only the finding may read Chapter~13 and Chapter~14 and stop.
Readers who want the mathematics may read Part~II first. Readers who want to
know what a proof assistant can and cannot certify may read Chapter~6 and
Chapter~9 and skip the rest.

One note on scope and honesty. This book is written by someone who is not a
number theorist, and who says so rather than hiding it. The analytic number
theory is explained from first principles wherever it is used; the claims about
the state of the literature are sourced and date-stamped; and where the audit
reached the limit of what could be checked without running the compiler, that
limit is marked in the text rather than papered over. The corrections log that
accompanies this book records, in the open, the places where the writer's
initial assumptions had to be revised in the course of the work --- including
several that would have made the book wrong.

% ---- How to use this book -----------------------------------------------
\chapter*{How to Use This Book}
\addcontentsline{toc}{chapter}{How to Use This Book}
\markboth{How to Use This Book}{How to Use This Book}

\thispagestyle{plain}

\begin{takeaway}{The shape of the book}
\begin{itemize}
\item \textbf{Part I} reports the claim and the artifact: what was said on
20 November 2026, and what was actually submitted.
\item \textbf{Part II} gives the mathematics. Chapter~3 builds the Riemann
zeta function from first principles and states the hypothesis precisely;
Chapter~4 sets out the four obstacles any method must get past; Chapter~5
reconstructs the method from the author's own timeline.
\item \textbf{Part III} is the audit. Chapter~6 explains what a proof assistant
certifies and what it does not. Chapter~7 traces the proof chain. Chapter~8
examines the load-bearing inputs. Chapter~9 reports the static scan.
\item \textbf{Part IV} places the claim in its context: the independent
compilations, the two-thirds theorem that preceded it, and the absence of any
independent review.
\item \textbf{Part V} gives the finding, the protocol that would settle it, and
what a proof would actually change.
\end{itemize}
\end{takeaway}

Three readings are possible, and each is coherent.

\emph{If you want to know what happened.} Read Chapters~1, 2, 10, 12 and 13.
That is the news, the artifact, the compilations, the state of the review, and
the finding.

\emph{If you want the mathematics.} Read Chapters~3, 4 and 5, then 8. These
assume no background beyond undergraduate analysis, and every result the
argument depends on is stated, sourced, and explained where the statement alone
is not enough.

\emph{If you want to check the book.} The audit in Part~III is written so that
it can be repeated. Chapter~7 gives the exact file paths and the exact theorem
names, so that the chain can be walked independently. Chapter~14 gives the
commands. Chapter~9 gives the scan so that the scan can be re-run and compared,
including the possibility that a reader's numbers differ.

\begin{mythbox}{Three things this book will not do}
\begin{itemize}
\item It will not declare the Riemann Hypothesis proved. The evidence assembled
here does not support that, and the author of the claim himself says he has
``just taken a step''.
\item It will not declare the claim false. Nothing found in this audit refutes
it, and the artifact is of a quality that makes refutation implausible on the
strength of a reading of the remark.
\item It will not treat a missing referee report as evidence of a defect. The
absence of review is a fact about the world, not a fact about the mathematics.
It is reported as the first and largest obstacle to a verdict.
\end{itemize}
\end{mythbox}

\vspace{8pt}
\noindent{\small\color{txtdim}\sffamily
Conventions. Mathematical statements are given in ordinary notation. Code is
set in a monospaced face and quoted exactly as it appears in the submitted
source, including its file paths. Where the book says that something was
``verified'', it means it was checked against the primary record named in the
Notes and References --- not that it was recalled. Where the book says
``unverified'', it means exactly that, and the sentence in question says what
would be needed to verify it.\par}
